\documentclass[10pt,aspectratio=169]{beamer}
\usepackage{amsfonts,amsmath,booktabs,graphicx}
\usetheme{sintef}
\usepackage{xeCJK}
\setsansfont{Times New Roman}
\usepackage{unicode-math}
\setmathfont{Latin Modern Math}
\setCJKsansfont{LXGW WenKai Screen}
\usefonttheme[onlymath]{serif}
\titlebackground*{assets/background-szu}
\graphicspath{{hair_parting_assets/}}
\setbeamertemplate{frametitle}{%
  \vspace*{-3.5ex}
  \begin{beamercolorbox}[leftskip=2cm]{frametitle}%
    \usebeamerfont{frametitle}\insertframetitle
  \end{beamercolorbox}
}

\title{带发缝头发的方向场建模}
\subtitle{从发缝定位到分区 PDE}
\author{Xie Xiang}
\date{2026/09/05}

\newcommand{\result}[1]{\textcolor{blue!70!black}{#1}}
\newcommand{\fail}[1]{\textcolor{red!70!black}{#1}}

\begin{document}
\maketitle

\begin{frame}{现在的问题：头发为什么会糊在一起？}
\Large
当前模型把整个头发区域当成一个连续的 PDE 域。

\vspace{0.8em}
\begin{center}
\textbf{同一邻域内持续平滑}
\quad$\Longrightarrow$\quad
\textbf{发缝两侧互相影响}
\quad$\Longrightarrow$\quad
\textbf{发丝跨缝、混在一起}
\end{center}

\vspace{1em}
\normalsize
即使二维发缝位置准确，如果三维求解仍允许跨岸传播，发缝也会被重新填满。

\vfill
\centering
\result{核心问题不是“发缝线不够细”，而是“求解域没有真正分开”。}
\end{frame}

\begin{frame}{这次准备怎么改？}
\begin{center}
\Large
精确发缝定位
$\longrightarrow$
左右分区
$\longrightarrow$
三维标签
$\longrightarrow$
分区 PDE
\end{center}

\vspace{1.4em}
\begin{columns}[T]
\column{0.30\textwidth}
\textbf{看得准}

从原图确定发缝位置。

\column{0.30\textwidth}
\textbf{分得开}

用 strand/depth 补充分区证据。

\column{0.30\textwidth}
\textbf{算得稳}

在 PDE 中切断跨区耦合。
\end{columns}

\vfill
\centering
这不是在结果上“画一道缝”，而是改变方向场的传播方式。
\end{frame}

\begin{frame}{数学模型只需要记住两件事}
\begin{columns}[T]
\column{0.46\textwidth}
\textbf{1. 头发方向贴着头皮走}

头皮是曲面，方向属于每个点的切平面。普通二维拉普拉斯不能直接描述曲率带来的标架变化。

\vspace{1em}
\textbf{2. 发丝方向没有正反}

$\theta$ 与 $\theta+\pi$ 表示同一条发流，因此使用二重角线场
\[
q_i=\exp(2\mathrm i\theta_i).
\]

\column{0.50\textwidth}
\textbf{Connection Laplace--Beltrami 能量}
\[
E(q)=\frac12\sum_{(i,j)}w_{ij}
\left|q_j-r_{ij}^{2}q_i\right|^2.
\]

\begin{itemize}
  \item $w_{ij}$：曲面上的余切权重
  \item $r_{ij}$：相邻切平面之间的平行移动
  \item 最小化 $E$：得到平滑的曲面发流
\end{itemize}
\end{columns}
\end{frame}

\begin{frame}{Step 1--3：先把发缝放到正确的曲面位置}
\begin{columns}[T]
\column{0.32\textwidth}
\centering
\includegraphics[width=\linewidth,height=0.54\textheight,keepaspectratio]{step01.png}

\textbf{Step 1}\quad 原图发缝反投影

\column{0.32\textwidth}
\centering
\includegraphics[width=\linewidth,height=0.54\textheight,keepaspectratio]{step02.png}

\textbf{Step 2}\quad 建立双岸与空带

\column{0.32\textwidth}
\centering
\includegraphics[width=\linewidth,height=0.54\textheight,keepaspectratio]{step03.png}

\textbf{Step 3}\quad 验证曲面算子
\end{columns}

\vfill
\centering
\result{结论：发缝位置、曲面拓扑和 Connection-LB 算子均可用。}
\end{frame}

\begin{frame}{Step 4--5：得到自然的两岸发流}
\begin{columns}[T]
\column{0.49\textwidth}
\centering
\includegraphics[width=\linewidth,height=0.57\textheight,keepaspectratio]{step04.png}

\textbf{Step 4：两岸张角}

最终采用 $60^\circ$，分流清楚且不过度张开。

\column{0.49\textwidth}
\centering
\includegraphics[width=\linewidth,height=0.57\textheight,keepaspectratio]{step05.png}

\textbf{Step 5：融合 strand map}

保留原图中的整体发型走向。
\end{columns}
\end{frame}

\begin{frame}{Step 6--8：局部结果是成立的}
\begin{columns}[T]
\column{0.49\textwidth}
\centering
\includegraphics[width=\linewidth,height=0.58\textheight,keepaspectratio]{step06.png}

沿曲面方向场积分 256 根根部轨迹。

\column{0.49\textwidth}
\centering
\includegraphics[width=\linewidth,height=0.58\textheight,keepaspectratio]{step08.png}

与外层发丝耦合后，局部发缝可见率为
\result{70.10\%}。
\end{columns}

\vfill
\centering
这说明根部方向场本身没有根本性错误。
\end{frame}

\begin{frame}{Step 9：扩展到 10k 后，发缝又被遮住}
\begin{columns}[T]
\column{0.58\textwidth}
\centering
\includegraphics[width=\linewidth,height=0.72\textheight,keepaspectratio]{step09.png}

\column{0.38\textwidth}
\Large
局部 256 根

\result{70.10\%}

\vspace{1.1em}
完整 10k

\fail{33.14\%}

\vspace{1.1em}
\normalsize
根部正确，并不代表完整发型中发缝可见。
\end{columns}
\end{frame}

\begin{frame}{Step 10：真正遮挡发缝的是中长段发丝}
\begin{columns}[T]
\column{0.58\textwidth}
\centering
\includegraphics[width=\linewidth,height=0.72\textheight,keepaspectratio]{step10a.png}

\column{0.38\textwidth}
在发缝投影区域内，共找到

\vspace{0.5em}
\centering
{\Huge\fail{1,289}}

\normalsize
条真实遮挡发丝。

\vspace{1em}
旧方法只选中 256 根，其中真正命中 20 根，召回率仅
\fail{1.55\%}。
\end{columns}
\end{frame}

\begin{frame}{为什么直接“把发丝移开”没有成功？}
\begin{columns}[T]
\column{0.56\textwidth}
\centering
\includegraphics[width=\linewidth,height=0.72\textheight,keepaspectratio]{step10c.png}

\column{0.40\textwidth}
\begin{itemize}
  \item 减少投影遮挡时，容易产生穿插和轮廓破坏
  \item 局部移动无法稳定处理 1,289 条长距离遮挡
  \item trust-region 候选全部被门禁拒绝
\end{itemize}

\vspace{1em}
\result{因此要把“左右归属”提前到场求解阶段。}
\end{columns}
\end{frame}

\begin{frame}{新的分区方法：位置证据 + 结构证据}
\begin{columns}[T]
\column{0.60\textwidth}
\centering
\includegraphics[width=\linewidth,height=0.74\textheight,keepaspectratio]{partition.png}

\column{0.36\textwidth}
\textbf{发缝位置}

RGB / Lab / DINO 给出高精度中心线。

\vspace{1em}
\textbf{分区结构}

strand map 判断方向突变，depth map 判断前后层次变化。

\vspace{1em}
\textbf{有效范围}

所有判断都限制在头发 seg 内。
\end{columns}
\end{frame}

\begin{frame}{分区 PDE：从算子层面阻止跨缝混合}
给每个网格点一个分区标签 $\ell_i$，只在同一区域内进行平滑：
\[
E_{\mathrm{part}}(q)=
\frac12\sum_{(i,j)}w_{ij}\,
\mathbf 1[\ell_i=\ell_j]
\left|q_j-r_{ij}^{2}q_i\right|^2
+E_{\mathrm{constraint}}.
\]

\vspace{0.8em}
\begin{columns}[T]
\column{0.46\textwidth}
\textbf{同一标签}

保留原来的曲面平滑与图像观测。

\column{0.46\textwidth}
\textbf{不同标签}

边权置零，不允许方向场跨过发缝传播。
\end{columns}

\vfill
\centering
\result{发缝从“软约束”变成 PDE 域的真实内部边界。}
\end{frame}

\begin{frame}{下一步只验证三件事}
\begin{enumerate}
  \setlength{\itemsep}{1.2em}
  \item \textbf{二维分区能否稳定提升到三维头皮}
  \item \textbf{分区 PDE 是否减少跨岸方向传播}
  \item \textbf{RK4 积分是否始终留在根部所属分区}
\end{enumerate}

\vspace{1.5em}
通过这三项后，再运行完整 10k 重建并比较发缝可见率、穿插率和轮廓变化。

\vfill
\footnotesize
可微投影占据率仍有价值，但放在分区 PDE 验证之后。
\end{frame}

\begin{frame}{结论}
\Large
\begin{enumerate}
  \setlength{\itemsep}{1.1em}
  \item Connection-LB 能得到合理的曲面根部方向场。
  \item 完整发型中的主要问题是跨岸耦合和中长段遮挡。
  \item 下一版模型应采用混合分区，并在 PDE 中切断跨区传播。
\end{enumerate}

\vfill
\centering
\Huge Questions?

\vspace{0.8em}
\tiny
参考：Knöppel et al., 2013；Crane et al., 2010；Dr.Hair, CVPR 2024。
\end{frame}

\end{document}
